% s-354-p354-m002; original PDF 354, 355, 356, 357, 358
\Aphorism{s-354-p354-m002}{405}
Man kann in einem doppelten Sinne von der ewigen Wiederkunft sprechen. Eine beliebige, erfüllte Zeitstrecke kann in der absoluten Zeit beliebig oft hintereinander abgespielt werden, ohne dass diese Wiederholung ins empirische Bewusstsein der die Zeitstrecke erlebenden Wesen fällt. Dies gilt unabhängig vom Inhalt der Strecke, als ein aus dem Begriff der Zeit folgender Satz \Foreign{a priori.} Andererseits wäre es möglich, dass unser empirischer Zeitinhalt als Ganzes eine geschlossene Linie darstellte; das ist der Sinn, in welchem Nietzsche die uralte Vorstellung einer ewigen Wiederkehr aller Dinge als poetische Hypothese wieder aufgenommen hat. Diese Wiederkehr ist aber eine Aussage \Foreign{a posteriori,} geschöpft aus dem Inhalt unserer speciellen Zeitlinie, des Ensembles oder Continuums von Weltzuständen, das wir bilden und erfahren: während jene erste Art Wiederkehr auch ohne die zweite besteht und, wie gesagt, im Wesen des zeitlichen Erlebens begründet ist. — Die zweite, inhaltlich zu begründende Wiederkunftslehre besagt also, dass es zwar unbegrenzt viele Weltzustände giebt, die aber so angeordnet sind, dass man von irgend einem derselben mit der Zeit fortschreitend wieder zu ihm zurückgelangt. Die Zeitlinie, der wir angehören, wäre demnach, als geometrische Linie gedacht, etwa ein Kreis, während sie dem naiven Menschen als Gerade erscheint. Also derselbe Fall, wie bei der Erde, deren Oberfläche dem ersten Anschein nach eine Ebene ist; die Krümmung der Kugel nämlich ist zu gering, um sich innerhalb eines Flächenstücks, das der Mensch mit seiner normalen Thätigkeit umspannt, der roheren Beobachtung zu verrathen. Man muss, um die Krümmung zu constatiren, entweder die Beobachtung verfeinern oder ihr Gebiet erweitern; das Äusserste in der einen Richtung leistet die Geodäsie, die den Winkel der Verticalen schon innerhalb eines Gärtchens bestimmen kann; das Extrem in der anderen wäre etwa eine Weltumsegelung. Nun, ebenso steht es mit der ewigen Wiederkunft — auch dieser Kreis muss eine stärkere oder schwächere Krümmung, d. h. seinen kleineren oder grösseren Radius haben; gross oder klein im Verhältniss zu einer bekannten Zeiteinheit, vielleicht einem Sonnenjahr oder einem Menschenleben oder einem Jahrhundert. Dass auch hier die „Krümmung“ nicht beträchtlich sein kann, zeigt sich darin, dass sie eben dem naiven Bewusstsein ganz entgeht, selbst wenn es sich mit Historie bewaffnet; in der That deutet bis jetzt nichts darauf hin, dass wir uns einem geschichtlich überlieferten Zeitraum wiederum \Foreign{a parte ante} näherten. Trotzdem ist Nietzsches Hypothese wohl denkbar, unter der plausiblen Annahme, dass im historischen Bewusstsein Discontinuitäten stattgefunden haben, veranlasst etwa durch Erdrevolutionen oder Katastrophen kosmischen Charakters, sodass dem Menschen der\Sperr{ historische Rückblick bis zu sich selber }versperrt wäre\Sperr{.} Dass wir von einer solchen Zwischenzeit historischen Unbewusstseins herkommen, wissen wir: sie liegt noch gar nicht lange, der Schätzung nach zehn Jahrtausende, zurück. Dass wir \Sperr{möglicherweise} einer gleichen, oder vielmehr \Sperr{derselben} Unterbrechung des geschichtlichen Zusammenhanges der Menschheit zusteuern, wird Niemand leugnen können. Ohne Annahme einer solchen Unterbrechung würde allerdings die Krümmung des Kreises sehr klein, sein Radius sehr gross vorausgesetzt werden müssen, sodass die Nachwirkung eines Ereignisses oder Zeitraums durch die ganze Peripherie des Zeitkreises hin bis zu sich selber (die Erinnerung des Ereignisses an sich selber, kann man ebensogut sagen) \Sperr{zu schwach} wäre, um dem Bewusstsein aufzufallen. — Verlassen wir aber einmal das wirklich Gegebene und gestatten uns einen Phantasieeinblick ins Mögliche, so hindert nichts, geschlossene Zeitlinien zu construiren, deren Peripherie in \Sperr{mässig grossem} Verhältniss zu einer der genannten menschlichen oder natürlichen Einheiten steht, und auf denen zweitens keine Unterbrechung des historischen Bewusstseins eintritt. Die ewige Wiederkunft wäre dann keine Hypothese mehr, sondern eine Erfahrungsthatsache, der sich die einigermassen gesteigerte menschliche Besonnenheit nicht entziehen könnte; beispielsweise würde, die Peripherie der Zeitlinie zu tausend Jahren gerechnet, der Mensch zum Vorfahren in 29ter Generation seinen eigenen Sohn, in 30ter sich selbst haben, oder auch, er würde den Schluss ziehen können, dass sein Abkömmling in 29ter Generation sein eigener Vater, in 30ter er selbst sein wird. Der gesammte Zeitinhalt würde einem historisch genügend fernsichtigen Bewusstsein bekannt sein, also, in Form der Vergangenheit, auch die ganze Zukunft, woraus sich für die Motivation des willkürlichen Handelns Verhältnisse ergeben müssen, die uns so ziemlich unvorstellbar sind — unser menschliches Wollen und Wirken beruht ja auf Unkenntniss der Zukunft! — Vergrössert man die Krümmung des Zeitkreises noch mehr, so bedarf es nicht mehr der historischen Teleskopie, um das periodische Wiederkehren aller Dinge als Bewusstseinsobject in deutlicher Sehweite zu haben; die Erfahrung weniger Generationen, eines einzelnen Individuums würde ausreichen — vielleicht (wenn man den Kreis immer kleiner werden lässt) die einer Stunde, die einer Minute! Einer Stunde oder Minute natürlich nicht im astronomischen Sinne; denn wenn der ganze Weltkreisprocess eine Stunde lang ist und dann in sich zurückläuft, so kann unsere Erde weder ihre tägliche Drehung noch ihre jährliche Bahn vollenden. Es ist aber an die physiologische Bedeutung dieser Zeitmasse zu denken, die von der Organisation des erlebenden Wesens abhängt; als Secunde wäre irgend ein Vielfaches der Dauer eines psychischen Elementarvorgangs, sei es Muskelcontraction oder Nervenreaction, zu definiren. — Wie allerdings einem Wesen zu Muthe sein mag, dessen Zeitperipherie eine Stunde lang ist, davon giebt uns die ausschweifendste Phantasie keine rechte Vorstellung. Wenn wir räumlich in einem Kreise herumgeführt werden, so stellt sich der Schwindel ein, desto stärker, je stärker der Kreis gekrümmt ist. An nichts Geringerem werden wohl die armen Gehirne zu leiden haben, die mit Bewusstsein ihren zeitlichen Rundlauf im engsten Kreise beginnen und vollenden und ewig wieder beginnen müssen. Dass uns kein solches Loos fiel, dass uns als Erkenntniss-Subjecten ein Kreis kleiner oder verschwindender Krümmung zugeordnet ist, dafür dürfen wir uns schon — bei uns selbst — bedanken!
\par

